\documentclass[a4paper,twoside]{article}

\usepackage{morewrites}

\usepackage[svgnames,dvipsnames,x11names]{xcolor}
\usepackage{graphicx}
\usepackage[english]{babel}
\usepackage[colorlinks=true, linkcolor=NavyBlue, urlcolor=NavyBlue, citecolor=NavyBlue]{hyperref}
\usepackage[a4paper]{geometry}
\usepackage{ts-fonts}
\usepackage{booktabs}
\usepackage{tabularx}
\usepackage{amsmath}
\usepackage{float}
\usepackage{seqsplit}
\newcolumntype{Y}{>{\centering\arraybackslash}X}
\newcolumntype{Z}{>{\raggedleft\arraybackslash}X}
\usepackage{ts-code}

\usepackage{lastpage}
\usepackage{refcount}
\makeatletter
\def\ps@plain{
  \let\@oddhead\@empty
  \let\@evenhead\@empty
  \def\@oddfoot{
    \hfil
    \ifnum\getpagerefnumber{LastPage}>1\relax
      \thepage
    \fi
    \hfil}
  \let\@evenfoot\@oddfoot
}
\makeatother

\title{Integration with MkDocs}
\author{}\date{}

\begin{document}

\maketitle

\thispagestyle{plain}
TeXSmith can be seamlessly integrated with MkDocs, a popular static site generator for project documentation. This allows you to leverage MkDocs' powerful features while utilizing TeXSmith for document generation.

The integration is achieved through the \texttt{mkdocs-\allowbreak{}texsmith} plugin, which processes TeXSmith documents during the MkDocs build process.

\section{Configuration}\label{configuration}

To enable TeXSmith in your MkDocs project, you need to add the \texttt{mkdocs-\allowbreak{}texsmith} plugin to your \texttt{mkdocs.yml} configuration file:

\begin{code}{yaml}{}{}
\begin{Verbatim}[commandchars=\\\{\},breaklines, breakanywhere, commandchars=\\\{\}]
\PY{n+nt}{plugins}\PY{p}{:}
\PY{+w}{  }\PY{p+pIndicator}{\PYZhy{}}\PY{+w}{ }\PY{l+lScalar+lScalarPlain}{texsmith}
\end{Verbatim}
\end{code}
You can configure additional options for the TeXSmith plugin as needed:

\begin{center}
\begin{tabularx}{\linewidth}{>{\raggedright\arraybackslash}X>{\raggedright\arraybackslash}X>{\raggedright\arraybackslash}X}
\toprule
\textbf{Option} & \textbf{Description} & \textbf{Default} \\

\midrule
\texttt{template} & Template to use for rendering the site & \texttt{book} \\
\texttt{build\_dir} & Directory where TeXSmith outputs are stored & \texttt{site} \\

\bottomrule
\end{tabularx}
\end{center}

\subsection{Multiple documents}\label{multiple-documents}

You can either generate a single document from your MkDocs site or multiple documents from different sections.

\begin{code}{yaml}{}{}
\begin{Verbatim}[commandchars=\\\{\},breaklines, breakanywhere, commandchars=\\\{\}]
\PY{n+nt}{plugins}\PY{p}{:}
\PY{+w}{    }\PY{p+pIndicator}{\PYZhy{}}\PY{+w}{ }\PY{n+nt}{texsmith}\PY{p}{:}
\PY{+w}{      }\PY{n+nt}{books}\PY{p}{:}
\PY{+w}{        }\PY{p+pIndicator}{\PYZhy{}}\PY{+w}{ }\PY{n+nt}{template}\PY{p}{:}\PY{+w}{ }\PY{l+lScalar+lScalarPlain}{book}
\PY{+w}{          }\PY{n+nt}{folder}\PY{p}{:}\PY{+w}{ }\PY{l+lScalar+lScalarPlain}{foolists}
\PY{+w}{          }\PY{n+nt}{root}\PY{p}{:}\PY{+w}{ }\PY{l+s}{\PYZdq{}}\PY{l+s}{foo}\PY{l+s}{\PYZdq{}}
\PY{+w}{          }\PY{n+nt}{base\PYZus{}level}\PY{p}{:}\PY{+w}{ }\PY{l+lScalar+lScalarPlain}{\PYZhy{}1}
\PY{+w}{        }\PY{p+pIndicator}{\PYZhy{}}\PY{+w}{ }\PY{n+nt}{template}\PY{p}{:}\PY{+w}{ }\PY{l+lScalar+lScalarPlain}{article}
\PY{+w}{          }\PY{n+nt}{folder}\PY{p}{:}\PY{+w}{ }\PY{l+lScalar+lScalarPlain}{bariers}
\PY{+w}{          }\PY{n+nt}{root}\PY{p}{:}\PY{+w}{ }\PY{l+s}{\PYZdq{}}\PY{l+s}{bar}\PY{l+s}{\PYZdq{}}
\end{Verbatim}
\end{code}
\section{Serve}\label{serve}

During development with \texttt{mkdocs serve}, the TeXSmith plugin can fetch assets from the web (e.g. images, citations) and compile PDF snippets on the fly. This allows for a smooth writing experience with instant feedback.

\section{Build}\label{build}

When you run \texttt{mkdocs build}, the TeXSmith plugin processes all TeXSmith documents in your project, generating the corresponding PDFs and integrating them into the final site output. By default the output directory is \texttt{press/}.

You can tell TeXSmith to build the PDF during the process with the \texttt{build} option in your mkdocs.yml:

\begin{code}{yaml}{}{}
\begin{Verbatim}[commandchars=\\\{\},breaklines, breakanywhere, commandchars=\\\{\}]
\PY{n+nt}{plugins}\PY{p}{:}
\PY{+w}{  }\PY{p+pIndicator}{\PYZhy{}}\PY{+w}{ }\PY{n+nt}{texsmith}\PY{p}{:}
\PY{+w}{      }\PY{n+nt}{build}\PY{p}{:}\PY{+w}{ }\PY{l+lScalar+lScalarPlain}{true}
\end{Verbatim}
\end{code}

\end{document}
